% Development supplement, 11 October 2026.
% These are analytic transport lemmas, not a proof of the new local moment.
\section{Uniform deletion and common Fourier transport}
\label{sec:uniform-transport}

The following estimates resolve two bookkeeping questions in the proposed
argument. They neither identify a transformed character nor supply a
missing mean-value estimate. The field throughout is the fixed Eisenstein
field. Constants may depend on the fixed field, the indicated compact
supports, and the requested error exponent.

\begin{proposition}[Uniform finite Euler deletion]
Fix $B\geq 0$, $\sigma_0>0$, an integer $j\geq0$, and $\epsilon>0$.
Let $Z\geq2$, let $\mathfrak q$ be a squarefree integral ideal with
$N\mathfrak q\leq Z^B$, and choose complex numbers
$|a_{\mathfrak p}|\leq1$. Put
\[
 E_{\mathfrak q}(s)=\prod_{\mathfrak p\mid\mathfrak q}
       (1-a_{\mathfrak p}(N\mathfrak p)^{-s}).
\]
Then $E_{\mathfrak q}$ is nonzero on $\Re s>0$, and, uniformly in
$\mathfrak q$, the coefficients, and $\Im s$,
\[
 |E_{\mathfrak q}^{(j)}(s)|+
 |(E_{\mathfrak q}^{-1})^{(j)}(s)|
 \ll_{B,\sigma_0,j,\epsilon} Z^\epsilon
 \qquad(\Re s\geq\sigma_0).
\]
\end{proposition}
\begin{proof}
Every factor is nonzero because
$|a_{\mathfrak p}(N\mathfrak p)^{-s}|<1$ in the positive half-plane.
For $\Re s\geq\sigma_0/2$, both the product and its reciprocal are
bounded by
\[
 \prod_{\mathfrak p\mid\mathfrak q}
       (1-(N\mathfrak p)^{-\sigma_0/2})^{-1}.
\]
Given $b>0$, the factor at every sufficiently large prime norm is at
most $(N\mathfrak p)^b$: its logarithm is
$O((N\mathfrak p)^{-\sigma_0/2})=o(\log N\mathfrak p)$.
The finitely many smaller primes contribute a constant, independently
of the chosen ideal and coefficients. The product is consequently
$O((N\mathfrak q)^b)$. If $B>0$, choose $bB\leq\epsilon$; if $B=0$
the product is empty. Cauchy's derivative estimate on disks of radius
$\sigma_0/2$ proves the claim for every fixed $j$.
\end{proof}

Thus, once the exact identity
$L_{\rm masked}(s)=L_{\rm primitive}(s)E_{\mathfrak q}(s)$ is
established, extra polynomial-size holes preserve its positive-real
zeros and cost an arbitrarily small power in a contour argument staying
in a fixed positive half-plane. For a nonprincipal character the
identity also applies at $s=1$. A principal pole must still be treated
separately. This argument places no small-conductor hypothesis on the
\emph{masked presentation}. It applies the analytic estimate to the
primitive presentation and restores the finite factors afterward.
For the raw row family one still has to prove the primitive-conductor
bound and the exact coefficient identity; these are not consequences
of a product estimate.

For a prime polynomial with relative support $[c,C]\subset(0,\infty)$,
bounded coefficient modulus, and central normalization $P^{-1/2}$,
deleting primes dividing $\mathfrak q$ changes it by
\[
 O(P^{-1/2}\log(2N\mathfrak q))\qquad(P\geq1).
\]
Indeed there are at most $\log(N\mathfrak q)/\log2$ such primes.
Fixed powers of $\log N\mathfrak p$ add only powers of $\log Z$ on
polynomial-size scales. Real powers on the normalized support have
bounded modulus, and imaginary powers have modulus one. This error is
$O_\epsilon(Z^\epsilon)$ and is compatible with a squared pointwise
bound $P^\delta Z^\epsilon$ whenever $\delta\geq0$, after dividing
the error allowance between the original estimate and this correction.
No deletion of an arithmetic mask by positivity is involved.

\begin{proposition}[Fourier transport with labels inside the norm]
Fix a compact box $K\subset(0,\infty)^d$ and $J\geq0$. Suppose that
a family of smooth profiles $H_A$ is supported on one larger compact
box and has uniformly bounded Euler derivatives through order $2r$,
where $2r>J+d$. Then it has representations
\[
 H_A(y)=\int_{\mathbb R^d}\prod_{k=1}^d y_k^{it_k}\,d\mu_A(t),
 \qquad
 \int_{\mathbb R^d}(1+|t|)^J\,d|\mu_A|(t)\leq C_J,
\]
uniformly in $A$. The measure is chosen before any arithmetic label
or row is evaluated.

Let $\Lambda$ be a finite set of retained arithmetic labels, and let
$F_t=(F_t(\lambda,h))$ be a measurable family in
$\mathcal H=\bigoplus_{\lambda\in\Lambda}\ell^2(R_\lambda)$.
If
\[
 \|F_t\|_{\mathcal H}\leq B(1+|t|)^J,
\]
then
\[
 \left\|\int F_t\,d\mu_A(t)\right\|_{\mathcal H}
       \leq C_J B.
\]
In particular there is no factor $|\Lambda|$ in this transport.
The same conclusion holds for weighted direct sums with fixed
nonnegative arithmetic weights.
\end{proposition}
\begin{proof}
Pass to logarithmic coordinates $x_k=\log y_k$ and apply Fourier
inversion. Integrating by parts with $(1-\Delta)^r$ bounds the
Fourier density by $C(1+|t|)^{-2r}$. Compact support bounds the
necessary $L^1$ derivative norms uniformly. The asserted weighted
integrability follows because $2r>J+d$.

The second assertion is the triangle inequality for the Bochner
integral in the Hilbert space $\mathcal H$:
\[
 \left\|\int F_t\,d\mu_A(t)\right\|_{\mathcal H}
 \leq\int\|F_t\|_{\mathcal H}\,d|\mu_A|(t)
 \leq B\int(1+|t|)^J\,d|\mu_A|(t).
\]
The assumption supplies integrability. Alternatively, all row sums
may first be truncated, with the same estimate, followed by a limit.
\end{proof}

The required uniform derivative assertion holds for a kernel
$K(A\prod_k y_k^{e_k})$ if its Euler derivatives are bounded on
$(0,\infty)$ and the exponents $e_k$ are fixed. Explicitly,
\[
 D^\beta K\left(A\prod_k y_k^{e_k}\right)
   =\left(\prod_k e_k^{\beta_k}\right)
     (x\partial_x)^{|\beta|}K(x)
       \bigg|_{x=A\prod_k y_k^{e_k}}.
\]
Cutoffs on the fixed box and normalized norm powers introduce only
fixed seminorms and polynomials in their independent heights.
Relations among arithmetic norms restrict evaluation points inside
the box; they do not redefine $\mu_A$ for each label. Discontinuous
coprimality, valuation, or zero-bin indicators remain in the arithmetic
vector $F_t$ and are never differentiated.

The substantive premise is a bound on the \emph{joint} vector $F_t$.
Separate per-label bounds cannot be substituted for it without their
label count. Conversely, once an exact reconstruction bounds the joint
sum over $(\lambda,h)$, the Fourier step does not charge that label
count again. The preceding statement justifies this last passage but
does not prove the reconstruction or the required pointwise bounds.

There is also a useful version when an outer arithmetic label changes
the parameter $A$. The integration-by-parts proof gives a single
pointwise majorant for the Fourier densities,
$|f_A(t)|\leq C(1+|t|)^{-2r}$. Choose the positive measure
$d\mu(t)=C(1+|t|)^{-2r}\,dt$ before the labels. Write
$f_{A(\lambda)}(t)\,dt=a_\lambda(t)\,d\mu(t)$, where
$|a_\lambda(t)|\leq1$. For each $t$, multiplication by
$a_\lambda(t)$ is a diagonal contraction on the joint Hilbert space.
The same bound therefore follows with this one common dominating
measure; equality of the label-dependent Fourier densities is not
needed. This is used only for labels fixed outside the sum being
Poisson-transformed. It does not permit a discontinuous row-dependent
factor to be inserted into that sum before Poisson summation.

\paragraph{Quantifier order for height losses.}
If a density input is eventually bounded by
$C(1+T)^J H^\rho$, its threshold must be independent of $T$ before
one substitutes $T=Z^\tau$. An assertion proved only for each fixed
$T$ with an unspecified threshold $H_0(T)$ is insufficient. With a
uniform threshold, bounded $H$ are filled using the lattice count and
$(1+T)^J\geq1$, without a new height dependence. Fix the zero buffer,
the finitely many profiles, and their analytic height exponents first;
then choose $\tau$ so that the total $\tau J$ lies below the reserved
power saving. Frequencies outside the retained range require the
common Fourier tail estimate. The pointwise and density results must
be checked to have precisely this uniform order; it is not supplied
by the Fourier identity alone.
